\section*{Capítulo \ref{ch:g10:pressure} --- Pressão: do esporte ao mergulho}

\begin{solution}{exo:g10:pressure:1}
Face grande: $S = \qty{0.15}{m^2}$, $P = 600/0.15 = \qty{4.0}{kPa}$.
Face pequena: $S = \qty{0.060}{m^2}$, $P = \qty{10}{kPa}$. Na areia, a
face grande: \omterm{def:g10:pressure:pressure}{pressão} menor, menos afundamento.
\end{solution}

\begin{solution}{exo:g10:pressure:2}
$\qty{0.25}{MPa} = \qty{250}{kPa} = \qty{2.5e5}{Pa}$.
$3.0/\num{2.0e-4} = \qty{15}{kPa}$ vence
$60/0.050 = \qty{1.2}{kPa}$: a \omterm{def:g10:inertia:force}{força} pequena pressiona mais forte.
\end{solution}

\begin{solution}{exo:g10:pressure:3}
$\rho g h = 1000 \times 9.81 \times 3.0 \approx \qty{29}{kPa}$; absoluta
$P = 101 + 29 = \qty{130}{kPa}$.
\end{solution}

\begin{solution}{exo:g10:pressure:4}
$V_2 = 100 \times 4.0 / 250 = \qty{1.6}{L}$; a \qty{50}{kPa},
$V_2 = 400/50 = \qty{8.0}{L}$.
\end{solution}

\begin{solution}{exo:g10:pressure:5}
$F = 65 \times 9.81 \approx \qty{638}{N}$. Lâmina:
$638/\num{1.2e-3} \approx \qty{5.3e5}{Pa} = \qty{530}{kPa}$. Solas:
$638/0.036 \approx \qty{18}{kPa}$. Razão $\approx 30$: a lâmina morde o
gelo, as solas não.
\end{solution}

\begin{solution}{exo:g10:pressure:6}
Polegar: $20/\num{1.2e-4} \approx \qty{1.7e5}{Pa} = \qty{170}{kPa}$.
Parede: $20/\num{1.0e-8} = \qty{2.0e9}{Pa}$. Os mesmos \qty{20}{N} agem
sobre uma área $12\,000$ vezes menor na ponta, de modo que só ali a
\omterm{def:g10:pressure:pressure}{pressão} supera o que o material aguenta.
\end{solution}

\begin{solution}{exo:g10:pressure:7}
$\Delta P = \rho g h = 1025 \times 9.81 \times 30 \approx
\qty{3.0e5}{Pa}$, logo
$F = \num{3.0e5} \times 0.50 \approx \qty{1.5e5}{N}$ --- o \omterm{def:g10:universal-gravitation:weight}{peso} de
$F/g \approx \qty{1.5e4}{kg}$, quinze toneladas. Tripulação nenhuma abre
isso para fora contra o mar.
\end{solution}

\begin{solution}{exo:g10:pressure:8}
$\rho = \dfrac{P}{g h} = \dfrac{\num{66.2e3}}{9.81 \times 5.00}
\approx \qty{1350}{kg/m^3}$. Não é água doce ($\qty{1000}{kg/m^3}$): é um
líquido bem mais denso, por exemplo uma salmoura concentrada.
\end{solution}

\begin{solution}{exo:g10:pressure:9}
$h = P_0/(\rho g)$: água, $\num{1.01e5}/9810 \approx \qty{10.3}{m}$;
mercúrio, $\num{1.01e5}/(13600 \times 9.81) \approx \qty{0.76}{m}$. Um
tubo de \qty{76}{cm} cabe numa mesa; um barômetro de água de \qty{10}{m}
precisa de uma caixa de escada --- a densidade do mercúrio venceu.
\end{solution}

\begin{solution}{exo:g10:pressure:10}
$V = 6000/240 = \qty{25}{mL}$; $P = 6000/75 = \qty{80}{kPa}$. $P$ em
função de $V$ é uma hipérbole; $P$ em função de $1/V$ é uma reta pela
origem, de coeficiente angular \qty{6000}{kPa.mL}.
\end{solution}

\begin{solution}{exo:g10:pressure:11}
A \qty{20}{m}: $P = 101 + 2 \times 101 = \qty{303}{kPa}$, três
atmosferas. Na superfície,
$V = 0.50 \times 303/101 = \qty{1.5}{cm^3}$: a bolha triplica.
\end{solution}

\begin{solution}{exo:g10:pressure:12}
\emph{1.} $F = \num{1.01e5} \times \num{1.2e-3} \approx \qty{121}{N}$,
sustentando $m = 121/9.81 \approx \qty{12}{kg}$.

\emph{2.} $\Delta P = 101 - 20 = \qty{81}{kPa}$:
$F = \num{8.1e4} \times \num{1.2e-3} \approx \qty{97}{N}$, cerca de
\qty{9.9}{kg}.
\end{solution}

\begin{solution}{exo:g10:pressure:13}
Boyle: $12 \times (200 - 50) = \qty{1800}{L}$ de ar à \omterm{def:g10:pressure:pressure}{pressão} da
superfície são utilizáveis. A \qty{30}{m} cada respiração é tomada a
\qty{4}{bar}, de modo que \qty{15}{L/min} à \omterm{def:g10:pressure:pressure}{pressão} ambiente custam
$15 \times 4 = \qty{60}{L/min}$ de ar da superfície:
$1800/60 = \qty{30}{min}$. A \qty{10}{m} ($\qty{2}{bar}$):
\qty{30}{L/min}, logo \qty{60}{min}. A profundidade se paga em ar.
\end{solution}

\begin{solution}{exo:g10:pressure:14}
$P = P_0 \dfrac{V_0}{V} = 101 \times \dfrac{6.0}{1.5} = \qty{404}{kPa}$.
Com $\rho g \approx \qty{10.1}{kPa/m}$,
$h = (404 - 101)/10.1 \approx \qty{30}{m}$.
\end{solution}

\begin{solution}{exo:g10:pressure:15}
\emph{1.} \omterm{def:g10:universal-gravitation:weight}{Peso} $F = 1200 \times 9.81 \approx \qty{1.18e4}{N}$; \omterm{def:g10:pressure:pressure}{pressão}
$P = F/S = \num{1.18e4}/0.040 \approx \qty{2.9e5}{Pa}$; pistão pequeno:
$f = P s = \num{2.9e5} \times \num{2.0e-4} \approx \qty{59}{N}$ --- a
\omterm{def:g10:inertia:force}{força} fica dividida por $S/s = 200$.

\emph{2.} O volume de líquido se conserva:
$400 \times 2.0 = \qty{800}{cm^3}$ têm de vir do cilindro pequeno, que
percorre $800/2.0 = \qty{400}{cm} = \qty{4.0}{m}$. Trabalho:
$59 \times 4.0 \approx \qty{235}{J}$ de um lado,
$\num{1.18e4} \times 0.020 \approx \qty{235}{J}$ do outro --- o elevador
troca distância por \omterm{def:g10:inertia:force}{força}, nunca trabalho.
\end{solution}

\begin{solution}{pb:g10:pressure:1}
\textbf{1.} $\rho g = 1025 \times 9.81 \approx \qty{10.1}{kPa}$ por
\omterm{def:g10:orders-of-magnitude:unit}{metro}, de modo que \qty{10}{m} de água do mar acrescentam
$\approx \qty{101}{kPa}$: uma atmosfera a cada dez \omterm{def:g10:orders-of-magnitude:unit}{metros}, quase
exatamente.

\textbf{2.} $P = 101 + 10.1\,h$: \qty{202}{kPa} a \qty{10}{m},
\qty{302}{kPa} a \qty{20}{m}, \qty{403}{kPa} a \qty{30}{m} e
\qty{503}{kPa} a \qty{40}{m}.

\textbf{3.} $2.0$, $3.0$, $4.0$ e $5.0$ vezes $P_0$.

\textbf{4.} $\Delta P = 403 - 101 = \qty{302}{kPa}$ sobre
$\qty{0.015}{m^2}$:
$F = \num{3.02e5} \times 0.015 \approx \qty{4.5e3}{N}$ --- o \omterm{def:g10:universal-gravitation:weight}{peso} de
quase meia tonelada sobre o rosto.

\textbf{5.} $\Delta P = 10.1 \times 3.0 \approx \qty{30}{kPa}$;
$F = \num{3.0e4} \times \num{6.0e-5} \approx \qty{1.8}{N}$ --- um dedo
apertado contra o tímpano, já na profundidade de uma piscina. Compensar é
empurrar ar para o ouvido médio até que a \omterm{def:g10:pressure:pressure}{pressão} interna iguale a da
água.

\textbf{6.} $V = 6.0 \, P_0/P$: $6.0 \times 101/202 = \qty{3.0}{L}$;
$606/302 = \qty{2.0}{L}$; $606/403 = \qty{1.5}{L}$.

\textbf{7.} Cai à metade quando $P = 2P_0$, ou seja, a \qty{10}{m} --- já
no primeiro trecho do mergulho.

\textbf{8.} Não: $V = 606/P$ é uma hipérbole na \omterm{met:g10:pressure:dive}{pressão absoluta} (e é
$P$, não $h$, o que a \omterm{prop:g10:pressure:boyle}{lei de Boyle} enxerga); cada \qty{10}{m} a mais
retira menos volume do que os anteriores.

\textbf{9.} $V = \qty{1.5}{L}$ exige
$P = 101 \times 6.0/1.5 = \qty{404}{kPa}$, ou seja,
$h = 303/10.1 \approx \qty{30}{m}$ (\cref{exo:g10:pressure:14}).

\textbf{10.} O sangue: o plasma se desloca para os vasos do tórax, é
incompressível e ocupa o volume que o ar já não preenche.

\textbf{11.} À \omterm{def:g10:pressure:pressure}{pressão} ambiente, \qty{403}{kPa}. Por Boyle, o mesmo ar a
\qty{101}{kPa} ocuparia $4$ vezes o volume: ele é $4.0$ vezes mais denso
do que o ar da superfície.

\textbf{12.} $V = 6.0 \times 403/202 \approx \qty{12}{L}$ --- o dobro,
depois de apenas \qty{20}{m} de subida.

\textbf{13.} $V = 6.0 \times 403/101 \approx \qty{24}{L}$: quatro vezes a
capacidade pulmonar dele.

\textbf{14.} \emph{Nunca prenda a respiração na subida.} Com as vias
aéreas abertas, o ar em expansão sai pelo regulador, e o volume pulmonar
nunca ultrapassa \qty{6.0}{L}.

\textbf{15.} Os \qty{1.5}{L} de Lena a \qty{30}{m} \emph{são} os seus
\qty{6.0}{L} de superfície comprimidos: na subida eles se reexpandem até
exatamente \qty{6.0}{L}, e não mais. Já os \qty{6.0}{L} de Marco foram
\emph{tomados} a \qty{403}{kPa}: há \qty{24}{L} de ar de superfície
dentro deles.

\textbf{16.} $V = 1.0 \times 503/P$: \qty{1.25}{cm^3} a \qty{30}{m},
\qty{1.67}{cm^3} a \qty{20}{m}, \qty{2.49}{cm^3} a \qty{10}{m} e
\qty{4.98}{cm^3} na superfície.

\textbf{17.} Fatores por etapa: $503/403 = 1.25$, depois $1.33$, $1.49$
e $202/101 = 2.0$ nos últimos dez \omterm{def:g10:orders-of-magnitude:unit}{metros} --- o crescimento acelera à
medida que a superfície se aproxima.

\textbf{18.} $40/10 = \qty{4.0}{min}$. O último minuto: de \qty{10}{m}
até a superfície a \omterm{met:g10:pressure:dive}{pressão absoluta} cai à metade, a maior queda relativa
de toda a subida.

\textbf{19.} Uma subida rápida deixa o nitrogênio dissolvido efervescer
em bolhas dentro do sangue e das articulações (doença descompressiva);
subir devagar e fazer uma pausa perto de \qty{5}{m} mantém o gás
dissolvido tempo suficiente para que ele saia sem alarde pelos pulmões.

\textbf{20.} Apneísta: o ar dela apenas volta ao volume original ---
Boyle a protege. Mergulhador autônomo: o ar tomado em profundidade se
multiplica na subida --- respire, nunca prenda. E os últimos dez \omterm{def:g10:orders-of-magnitude:unit}{metros}
dobram todo volume aprisionado ($\times 2.0$, contra $\times 1.25$ lá
embaixo, a \qty{40}{m}): quanto mais perto da superfície, mais devagar
você deve se aproximar dela.
\end{solution}
